\documentclass[aspectratio=169]{beamer}
\usepackage[utf8]{inputenc}
\usepackage[spanish]{babel}
\usepackage{fontspec}
\usepackage{xcolor}
\usepackage{tikz}
\usepackage{booktabs}
\usepackage{array}
\usepackage{hyperref}

% --- Colors ---
\definecolor{bgdark}{HTML}{0D0D0D}
\definecolor{accent}{HTML}{6C63FF}
\definecolor{accentlight}{HTML}{A29BFE}
\definecolor{textwhite}{HTML}{F5F5F5}
\definecolor{textgray}{HTML}{B0B0B0}
\definecolor{success}{HTML}{00E676}
\definecolor{warning}{HTML}{FFD600}

% --- Theme ---
\usetheme{default}
\setbeamercolor{background canvas}{bg=bgdark}
\setbeamercolor{normal text}{fg=textwhite}
\setbeamercolor{frametitle}{fg=textwhite}
\setbeamercolor{itemize item}{fg=accent}
\setbeamercolor{itemize subitem}{fg=accentlight}
\setbeamerfont{frametitle}{size=\Large,series=\bfseries}
\setbeamertemplate{navigation symbols}{}
\setbeamertemplate{footline}{}
\setbeamertemplate{itemize item}{\textcolor{accent}{\rule{6pt}{6pt}}}
\setbeamertemplate{itemize subitem}{\textcolor{accentlight}{\small\textbullet}}

\begin{document}

% ============================================
% SLIDE 1: TITLE
% ============================================
\begin{frame}[plain]
\begin{center}
\vspace{1.2cm}
{\fontsize{14}{16}\selectfont\color{accentlight} AXIS --- OPEN WORLD $\vert$ BLOQUE 2}

\vspace{0.6cm}
{\fontsize{40}{44}\selectfont\bfseries\color{textwhite} VIBE CODING}

\vspace{0.3cm}
{\fontsize{18}{22}\selectfont\color{accent} Build Your Startup in 85 Minutes}

\vspace{1.2cm}
{\color{textgray}\large Breakout \& Lead PUCP}

\vspace{0.3cm}
{\color{textgray} 13 de Abril, 2026}
\end{center}
\end{frame}

% ============================================
% SLIDE 2: INTERACTIVE QUESTIONS
% ============================================
\begin{frame}[plain]
\begin{center}
\vspace{2cm}
{\fontsize{28}{32}\selectfont\color{textwhite}\bfseries ?`Qui\'enes tienen una idea\\[8pt] de negocio o proyecto?}

\vspace{2cm}
{\fontsize{18}{22}\selectfont\color{accent} Levanten la mano}
\end{center}
\end{frame}

% ============================================
% SLIDE 3: SECOND QUESTION
% ============================================
\begin{frame}[plain]
\begin{center}
\vspace{2cm}
{\fontsize{28}{32}\selectfont\color{textwhite}\bfseries ?`Y qui\'enes NO la han construido\\[8pt] porque no saben programar?}

\vspace{2cm}
{\fontsize{18}{22}\selectfont\color{warning} Esa excusa muere hoy.}
\end{center}
\end{frame}

% ============================================
% SLIDE 4: WHAT IS VIBE CODING
% ============================================
\begin{frame}{?`Qu\'e es Vibe Coding?}
\vspace{0.3cm}

{\large\color{accentlight} Construir software real describiendo lo que quieres.}

\vspace{0.5cm}

\begin{itemize}
\setlength\itemsep{0.6em}
\item T\'u describes tu idea en lenguaje natural
\item La IA escribe todo el c\'odigo
\item Ves resultados en tiempo real
\item Iteras en conversaci\'on: ``cambia esto'', ``agrega aquello''
\end{itemize}

\vspace{0.6cm}

\begin{center}
\begin{tikzpicture}
\node[rounded corners=8pt, fill=accent!15, text=textwhite, minimum width=13cm, minimum height=1.4cm, align=center] {
\large ``You fully give in to the vibes, embrace exponentials,\\and forget that the code even exists.'' --- Andrej Karpathy
};
\end{tikzpicture}
\end{center}
\end{frame}

% ============================================
% SLIDE 5: THE PROCESS
% ============================================
\begin{frame}{El Proceso}
\vspace{0.5cm}

\begin{center}
{\color{textgray} Lo que van a aprender hoy no es solo una herramienta. Es un \textbf{\color{accent}proceso}.}

\vspace{0.8cm}

\begin{tikzpicture}[
    box/.style={rounded corners=6pt, minimum width=3cm, minimum height=1.5cm, align=center, font=\normalsize},
    arrow/.style={->, very thick, color=accentlight}
]
\node[box, fill=accent!25, text=textwhite] (p1) at (0, 0) {\textbf{IDEA}\\[-1pt]\small Tu proyecto};
\node[box, fill=accent!35, text=textwhite] (p2) at (4, 0) {\textbf{PRD}\\[-1pt]\small via ChatGPT};
\node[box, fill=accent!50, text=textwhite] (p3) at (8, 0) {\textbf{BUILDER}\\[-1pt]\small Lovable};
\node[box, fill=accent!70, text=textwhite] (p4) at (12, 0) {\textbf{APP REAL}\\[-1pt]\small con IA};

\draw[arrow] (p1) -- (p2);
\draw[arrow] (p2) -- (p3);
\draw[arrow] (p3) -- (p4);
\end{tikzpicture}

\vspace{0.8cm}

{\color{textgray}\large Este proceso funciona con \textbf{cualquier} herramienta.\\Las herramientas cambian, el proceso queda.}
\end{center}
\end{frame}

% ============================================
% SLIDE 6: TOOLS
% ============================================
\begin{frame}{Herramientas de Vibe Coding}
\vspace{0.3cm}

\begin{columns}[T]
\column{0.48\textwidth}
{\large\color{accent}\textbf{Lovable}} {\small\color{success}(la de hoy)}
\begin{itemize}
\setlength\itemsep{0.2em}
\item Apps completas desde un prompt
\item Gratis, en el navegador
\item Deploy con un click
\end{itemize}

\vspace{0.4cm}
{\large\color{accentlight}\textbf{V0 by Vercel}}
\begin{itemize}
\setlength\itemsep{0.2em}
\item Interfaces y componentes
\item Excelente para UI
\end{itemize}

\column{0.48\textwidth}
{\large\color{accentlight}\textbf{Google AI Studio}}
\begin{itemize}
\setlength\itemsep{0.2em}
\item De Google, gratis con Gemini
\item Lo usaremos para el chatbot
\end{itemize}

\vspace{0.4cm}
{\large\color{accentlight}\textbf{bolt.new}}
\begin{itemize}
\setlength\itemsep{0.2em}
\item Similar a Lovable
\item Buena alternativa
\end{itemize}
\end{columns}

\vspace{0.5cm}
\begin{center}
{\color{textgray} Hoy usamos \textbf{\color{accent}Lovable} + \textbf{\color{accent}Google AI Studio}}
\end{center}
\end{frame}

% ============================================
% SLIDE 7: LIVE DEMO TRANSITION
% ============================================
\begin{frame}[plain]
\begin{center}
\vspace{2.5cm}
{\fontsize{32}{36}\selectfont\color{textwhite}\bfseries D\'enme una idea de startup.}

\vspace{0.8cm}
{\fontsize{20}{24}\selectfont\color{accent} Mientras m\'as loca, mejor.}

\vspace{1.5cm}
{\color{textgray}\large (Live Demo)}
\end{center}
\end{frame}

% ============================================
% SLIDE 8: SETUP - QR LOVABLE
% ============================================
\begin{frame}[plain]
\begin{center}
\vspace{1cm}
{\fontsize{24}{28}\selectfont\color{textwhite}\bfseries Abran dos cosas:}

\vspace{0.8cm}

\begin{columns}[c]
\column{0.48\textwidth}
\begin{center}
{\fontsize{18}{22}\selectfont\color{accent}\textbf{1. Lovable}}

\vspace{0.3cm}
{\large\color{textwhite} lovable.dev}

\vspace{0.2cm}
{\color{textgray} Creen cuenta con Google}

\vspace{0.3cm}
{\small\color{textgray} (QR en pantalla)}
\end{center}

\column{0.48\textwidth}
\begin{center}
{\fontsize{18}{22}\selectfont\color{accent}\textbf{2. Cualquier IA}}

\vspace{0.3cm}
{\large\color{textwhite} ChatGPT / Claude / Gemini}

\vspace{0.2cm}
{\color{textgray} En otra pesta\~na}

\vspace{0.3cm}
{\small\color{textgray} (la que prefieran)}
\end{center}
\end{columns}

\vspace{1cm}
{\color{warning}\large Sin laptop? Emp\'arejense con alguien.}
\end{center}
\end{frame}

% ============================================
% SLIDE 9: PRD CONCEPT
% ============================================
\begin{frame}{?`Qu\'e es un PRD?}
\vspace{0.3cm}

{\large\color{accentlight} Product Requirements Document}

\vspace{0.4cm}

Un documento que describe \textbf{exactamente} qu\'e quieres construir:
\begin{itemize}
\setlength\itemsep{0.4em}
\item Qu\'e p\'aginas tiene
\item Qu\'e hace cada secci\'on
\item Qu\'e texto muestra
\item Qu\'e colores y estilo usa
\end{itemize}

\vspace{0.5cm}

\begin{center}
\begin{tikzpicture}
\node[rounded corners=6pt, fill=accent!20, text=textwhite, minimum width=12cm, minimum height=1.2cm, align=center] {
\large Buen PRD = Buen resultado. \textbf{La IA lo escribe por ti.}
};
\end{tikzpicture}
\end{center}
\end{frame}

% ============================================
% SLIDE 10: PRD TEMPLATE
% ============================================
\begin{frame}{Template: Pega esto en tu IA favorita}
\vspace{0.1cm}

{\small
\begin{tikzpicture}
\node[rounded corners=8pt, fill=white!8, text=textwhite, minimum width=14.5cm, text width=14cm, align=left, inner sep=12pt] {
\texttt{I need you to create a detailed PRD for a landing page.}\\[4pt]
\texttt{Here's my project:}\\
\texttt{- Name: \textbf{\textcolor{warning}{[TU NOMBRE DE STARTUP]}}}\\
\texttt{- What it does: \textbf{\textcolor{warning}{[QUE HACE EN 1-2 FRASES]}}}\\
\texttt{- Target audience: \textbf{\textcolor{warning}{[PARA QUIEN ES]}}}\\
\texttt{- Color vibe: \textbf{\textcolor{warning}{[ej: ``dark and techy'']}}}\\ [4pt]
\texttt{Generate a PRD with: hero section, how it works (3 steps),}\\
\texttt{features (4-6), testimonials (3), pricing (3 tiers),}\\
\texttt{signup form, FAQ (5 questions), footer.}\\[4pt]
\texttt{For each section specify exact text, layout, and style.}\\
\texttt{Make it detailed enough to build without questions.}
};
\end{tikzpicture}
}

\vspace{0.3cm}
\begin{center}
{\color{accent}\textbf{Cambia las 4 variables en amarillo.} \color{textgray}Copia el PRD que te genera.}
\end{center}
\end{frame}

% ============================================
% SLIDE 11: BUILD TIME
% ============================================
\begin{frame}[plain]
\begin{center}
\vspace{1.5cm}
{\fontsize{36}{40}\selectfont\bfseries\color{accent} BUILD TIME}

\vspace{0.5cm}
{\fontsize{18}{22}\selectfont\color{textwhite} Pega tu PRD completo en Lovable.}

\vspace{0.3cm}
{\fontsize{18}{22}\selectfont\color{textwhite} Dale Enter. Espera la magia.}

\vspace{1.2cm}

{\color{textgray}\large Tips:}

\vspace{0.3cm}
{\color{accentlight} ``Cambia los colores a azul oscuro''}

\vspace{0.2cm}
{\color{accentlight} ``Traduce todo a espa\~nol''}

\vspace{0.2cm}
{\color{accentlight} ``Agrega animaciones al scroll''}

\vspace{0.2cm}
{\color{accentlight} ``Make the design more modern and clean''}
\end{center}
\end{frame}

% ============================================
% SLIDE 12: LEVEL UP - AI CHATBOT
% ============================================
\begin{frame}[plain]
\begin{center}
\vspace{2cm}
{\fontsize{28}{32}\selectfont\color{textwhite}\bfseries Ya tienen una landing page.}

\vspace{0.6cm}
{\fontsize{28}{32}\selectfont\color{accent}\bfseries Ahora le vamos a agregar\\[8pt] inteligencia artificial.}

\vspace{1.5cm}
{\fontsize{16}{20}\selectfont\color{textgray} Y es gratis.}
\end{center}
\end{frame}

% ============================================
% SLIDE 13: GEMINI API KEY
% ============================================
\begin{frame}{Paso 1: Obtener API Key de Gemini (gratis)}
\vspace{0.3cm}

\begin{center}
{\Large\color{accent}\textbf{aistudio.google.com/apikey}}

\vspace{0.3cm}
{\color{textgray} (QR en pantalla)}
\end{center}

\vspace{0.5cm}

\begin{enumerate}
\setlength\itemsep{0.5em}
\item Iniciar sesi\'on con tu cuenta de Google
\item Click en \textbf{``Create API Key''}
\item Seleccionar cualquier proyecto
\item \textbf{Copiar} la API key
\end{enumerate}

\vspace{0.5cm}

\begin{center}
\begin{tikzpicture}
\node[rounded corners=6pt, fill=warning!15, text=textwhite, minimum width=12cm, minimum height=0.9cm, align=center, font=\small] {
\textcolor{warning}{Nota:} API keys en frontend no es seguro para producci\'on. Para este workshop est\'a perfecto.
};
\end{tikzpicture}
\end{center}
\end{frame}

% ============================================
% SLIDE 14: SYSTEM PROMPT
% ============================================
\begin{frame}{Paso 2: System Prompt para tu chatbot}
\vspace{0.2cm}

{\small
\begin{tikzpicture}
\node[rounded corners=8pt, fill=white!8, text=textwhite, minimum width=14.5cm, text width=14cm, align=left, inner sep=12pt] {
\texttt{You are the AI assistant for \textbf{\textcolor{warning}{[STARTUP NAME]}}.}\\[3pt]
\texttt{\textbf{\textcolor{warning}{[STARTUP NAME]}} is \textbf{\textcolor{warning}{[QUE HACE EN 1-2 FRASES]}}.}\\[3pt]
\texttt{You help visitors understand the product,}\\
\texttt{answer questions about features and pricing,}\\
\texttt{and encourage them to sign up for early access.}\\[3pt]
\texttt{Be friendly, concise, and enthusiastic.}\\[3pt]
\texttt{If asked something you don't know, say:}\\
\texttt{``Great question! Sign up for early access}\\
\texttt{and our team will follow up with you.''}\\[3pt]
\texttt{Always respond in Spanish unless the user}\\
\texttt{writes in English.}
};
\end{tikzpicture}
}

\vspace{0.3cm}
\begin{center}
{\color{accent}\textbf{Cambia las variables en amarillo.} \color{textgray} Copia este texto.}
\end{center}
\end{frame}

% ============================================
% SLIDE 15: LOVABLE CHATBOT PROMPT
% ============================================
\begin{frame}{Paso 3: Pega esto en Lovable}
\vspace{0.2cm}

{\small
\begin{tikzpicture}
\node[rounded corners=8pt, fill=white!8, text=textwhite, minimum width=14.5cm, text width=14cm, align=left, inner sep=12pt] {
\texttt{Add an AI chatbot to my landing page:}\\[3pt]
\texttt{- Floating chat button, bottom-right, message icon}\\
\texttt{- Opens sleek chat window when clicked}\\
\texttt{- Welcome: ``Hola! Soy el asistente de}\\
\texttt{\ \ \textbf{\textcolor{warning}{[STARTUP NAME]}}. En que te puedo ayudar?''}\\
\texttt{- Google Gemini API (gemini-2.0-flash)}\\
\texttt{- API Key: \textbf{\textcolor{warning}{[TU API KEY]}}}\\
\texttt{- System instruction: \textbf{\textcolor{warning}{[TU SYSTEM PROMPT]}}}\\
\texttt{- Maintain chat history}\\
\texttt{- Match landing page colors}\\
\texttt{- ``Powered by AI'' badge at bottom}
};
\end{tikzpicture}
}

\vspace{0.3cm}
\begin{center}
{\color{accent}\textbf{Pega tus 3 valores} \color{textgray} y dale Enter. Tu landing ahora tiene IA.}
\end{center}
\end{frame}

% ============================================
% SLIDE 16: TEST YOUR CHATBOT
% ============================================
\begin{frame}[plain]
\begin{center}
\vspace{2cm}
{\fontsize{28}{32}\selectfont\color{textwhite}\bfseries Prueba tu chatbot}

\vspace{0.8cm}
{\fontsize{16}{20}\selectfont\color{accentlight} Preg\'untale sobre tu startup\\[8pt]
Preg\'untale los precios\\[8pt]
Preg\'untale algo que no sepa\\[8pt]
}

\vspace{1cm}
{\fontsize{14}{18}\selectfont\color{textgray} Si no funciona a la primera, dile a Lovable:\\
``Fix the chatbot, it's not connecting to the API''}
\end{center}
\end{frame}

% ============================================
% SLIDE 17: SHOWCASE TIME
% ============================================
\begin{frame}[plain]
\begin{center}
\vspace{2.5cm}
{\fontsize{36}{40}\selectfont\bfseries\color{accent} SHOWCASE}

\vspace{0.8cm}
{\fontsize{18}{22}\selectfont\color{textwhite} 3 voluntarios muestran lo que construyeron.}

\vspace{0.5cm}
{\fontsize{16}{20}\selectfont\color{textgray} 3 minutos cada uno.\\Landing + chatbot en vivo.}

\vspace{1cm}
{\fontsize{16}{20}\selectfont\color{warning} El mejor se lleva un libro.}
\end{center}
\end{frame}

% ============================================
% SLIDE 18: CLOSING
% ============================================
\begin{frame}[plain]
\begin{center}
\vspace{1.5cm}
{\fontsize{22}{26}\selectfont\color{textgray} Hace 85 minutos no sab\'ian programar.}

\vspace{0.6cm}
{\fontsize{28}{32}\selectfont\color{textwhite}\bfseries Ahora tienen una startup\\[8pt] con inteligencia artificial.}

\vspace{1cm}
{\fontsize{18}{22}\selectfont\color{accent} Las herramientas son gratis.\\[6pt] Las ideas son suyas.\\[6pt] Lo \'unico que falta es que empiecen.}
\end{center}
\end{frame}

% ============================================
% SLIDE 19: RESOURCES
% ============================================
\begin{frame}{Recursos}
\vspace{0.3cm}

\begin{columns}[T]
\column{0.48\textwidth}
{\color{accentlight}\textbf{Herramientas}}
\begin{itemize}
\setlength\itemsep{0.3em}
\item \textbf{Lovable:} lovable.dev
\item \textbf{V0:} v0.dev
\item \textbf{AI Studio:} aistudio.google.com
\item \textbf{bolt.new:} bolt.new
\end{itemize}

\vspace{0.4cm}
{\color{accentlight}\textbf{Deploy gratis}}
\begin{itemize}
\setlength\itemsep{0.3em}
\item \textbf{Vercel:} vercel.com
\item \textbf{Netlify:} netlify.com
\end{itemize}

\column{0.48\textwidth}
{\color{accentlight}\textbf{Para lanzar en serio}}
\begin{itemize}
\setlength\itemsep{0.3em}
\item \textbf{Dominio:} Namecheap (\$10/a\~no)
\item \textbf{Pagos:} Stripe
\item \textbf{Analytics:} Google Analytics
\end{itemize}

\vspace{0.4cm}
{\color{accentlight}\textbf{Breakout}}
\begin{itemize}
\setlength\itemsep{0.3em}
\item Pr\'oximamente: \textbf{Hackathon AI}
\item S\'igannos en redes: @breakout
\end{itemize}
\end{columns}

\vspace{0.5cm}
\begin{center}
{\color{warning}\textbf{Compartan lo que construyeron en IG/LinkedIn taggeando @breakout}}
\end{center}
\end{frame}

% ============================================
% SLIDE 20: FINAL
% ============================================
\begin{frame}[plain]
\begin{center}
\vspace{2.5cm}
{\fontsize{40}{44}\selectfont\bfseries\color{accent} Si puedes describirlo,}

\vspace{0.5cm}
{\fontsize{40}{44}\selectfont\bfseries\color{textwhite} puedes construirlo.}

\vspace{2cm}
{\color{textgray}\large Breakout \& Lead PUCP | AXIS 2026}
\end{center}
\end{frame}

\end{document}
